% chuong7.tex — phần mở rộng: đối chiếu hành vi với Luật An ninh mạng 2025 và Nghị định 330/2026
\section{Mở rộng: đối chiếu với Luật An ninh mạng 2025 và Nghị định 330/2026}\label{sec:7}

Chương 3 xét hành vi của taphoammo.net theo văn bản có hiệu lực tại thời điểm hành vi. Chương này trả lời một câu hỏi khác: nếu cùng các hành vi đó xảy ra khi Luật An ninh mạng 2025~\cite{luat-anm-2025} và Nghị định 330/2026/NĐ-CP~\cite{nd330-2026} đã có hiệu lực, thì khung pháp luật hiện hành gọi tên và xử phạt chúng ra sao. Mục đích là làm rõ khung hiện hành mà giải pháp ở chương 5 dựa vào, không phải áp ngược luật mới cho vụ việc. Các trích dẫn điều khoản lấy từ bản text của văn bản gốc trong thư mục \texttt{legal/}.

\subsection{Nguyên tắc áp dụng}\label{sec:7.1}

\begin{itemize}
\item \textbf{Hiệu lực.} Luật An ninh mạng 2025 có hiệu lực từ 1/7/2026; Luật An toàn thông tin mạng 2015 và Luật An ninh mạng 2018 hết hiệu lực từ cùng ngày (khoản 1, 2 Điều 44). Nghị định 330/2026 có hiệu lực từ 19/8/2026 (Điều 80).
\item \textbf{Không áp ngược.} Hành vi xảy ra trước 19/8/2026 áp dụng nghị định xử phạt có hiệu lực tại thời điểm vi phạm; chỉ áp dụng Nghị định 330/2026 nếu nghị định này không quy định trách nhiệm hoặc quy định nhẹ hơn (khoản 1 Điều 81).
\item \textbf{Hình sự loại trừ hành chính.} Hành vi có dấu hiệu tội phạm được chuyển hồ sơ truy cứu trách nhiệm hình sự (khoản 4 Điều 4). Ba bị can đã bị khởi tố theo Điều 288 Bộ luật Hình sự, nên các mức phạt dưới đây chỉ có ý nghĩa với người tham gia không bị truy cứu, hoặc với hành vi tương tự trong tương lai.
\item \textbf{Cách đọc mức phạt.} Ở Mục 1 đến Mục 5 Chương II (an ninh mạng), mức ghi trong điều là mức cho cá nhân, tổ chức chịu gấp hai lần; ở Mục 6 (bảo vệ dữ liệu cá nhân), mức ghi là mức cho tổ chức, cá nhân chịu một phần hai (khoản 1 Điều 7). Mức tối đa trong lĩnh vực an ninh mạng là 100 triệu đồng với cá nhân và 200 triệu đồng với tổ chức (khoản 3 Điều 7); với mua bán dữ liệu cá nhân là 10 lần khoản thu (điểm a khoản 4 Điều 7).
\item \textbf{Thời hiệu.} Thời hiệu xử phạt trong lĩnh vực an ninh mạng và bảo vệ dữ liệu cá nhân là 01 năm (khoản 1 Điều 3).
\end{itemize}

\subsection{Hành vi của nhà vận hành sàn}\label{sec:7.2}

Bảng dưới đây đặt các hành vi đã phân tích ở chương 3 cạnh điều cấm hoặc nghĩa vụ của Luật An ninh mạng 2025 và điều xử phạt của Nghị định 330/2026. Việc gắn hành vi với điều khoản là phân tích của nhóm.

{\small\setstretch{1.1}
\begin{xltabular}{\textwidth}{>{\hsize=1.0\hsize\raggedright\arraybackslash}X>{\hsize=0.95\hsize\raggedright\arraybackslash}X>{\hsize=0.85\hsize\raggedright\arraybackslash}X>{\hsize=1.2\hsize\raggedright\arraybackslash}X}
\caption{Hành vi của nhà vận hành sàn theo Luật An ninh mạng 2025 và Nghị định 330/2026}\label{tab:7.1}\\
\toprule
\textbf{Hành vi (mục ở chương 3)} & \textbf{Luật ANM 2025} & \textbf{NĐ 330/2026} & \textbf{Mức phạt tiền} \\
\midrule
\endfirsthead
\toprule
\textbf{Hành vi (mục ở chương 3)} & \textbf{Luật ANM 2025} & \textbf{NĐ 330/2026} & \textbf{Mức phạt tiền} \\
\midrule
\endhead
\bottomrule
\endlastfoot
Mua bán dữ liệu cá nhân, kể cả thông tin tài khoản số (3.1.1, 3.1.2) & Điểm h khoản 2 Điều 7; điểm g khoản 4 Điều 13 (và khoản 6 Điều 7 Luật BVDLCN 2025) & Khoản 1, 3 Điều 53 (Mục 6) & Từ 02 đến 10 lần khoản thu; nếu không có khoản thu: từ 70 triệu đến 3 tỷ đồng tùy số chủ thể dữ liệu và mức độ (tổ chức; cá nhân một phần hai) \\
Lập, vận hành, tiếp thị trực tuyến cho sàn trái phép (3.2.1, 3.2.3) & Điểm e khoản 4 Điều 13 & Điểm d khoản 3 Điều 14 & Cá nhân 20--30 triệu; tổ chức 40--60 triệu \\
Không kiểm duyệt, ngăn chặn, loại bỏ hành vi vi phạm trên hệ thống (3.2.1) & Điểm b khoản 2 Điều 15 & Điểm e khoản 2 Điều 15 & Cá nhân 20--30 triệu; tổ chức 40--60 triệu \\
Không lưu dữ liệu người dùng Việt Nam tại Việt Nam (3.2.2) & Khoản 3 Điều 25; khoản 7 Điều 41 & Điểm c khoản 2 Điều 29; điểm a khoản 1 Điều 33 & Điều 29: cá nhân 50--70 triệu, tổ chức 100--140 triệu. Điều 33: cá nhân 30--50 triệu, tổ chức 60--100 triệu \\
Cung cấp dịch vụ nhận tin nhắn, cuộc gọi để xác thực tài khoản (3.1.3) & Điểm g khoản 4 Điều 13 & Điểm c khoản 2 Điều 14, chỉ khi là tài khoản có chức năng giao dịch tài chính & Cá nhân 10--20 triệu; tổ chức 20--40 triệu \\
Trung gian thanh toán cho giao dịch hàng cấm (3.2.4) & Không có điều tương ứng trực tiếp & Điểm b khoản 2 Điều 14 (dùng tài khoản số làm công cụ trung gian nhận, chuyển tiền khi không được cấp phép); nếu là hành vi cung ứng dịch vụ trung gian thanh toán thì áp dụng nghị định chuyên ngành (khoản 3 Điều 1): khoản 6 Điều 31 Nghị định 340/2025/NĐ-CP~\cite{nd340-2025} & Điều 14: cá nhân 10--20 triệu; tổ chức 20--40 triệu. Nghị định 340/2025: cá nhân 150--250 triệu; tổ chức 300--500 triệu \\
\end{xltabular}}

Kèm theo phạt tiền, Nghị định 330/2026 quy định các biện pháp như: buộc hủy, xóa toàn bộ dữ liệu cá nhân đã mua bán, buộc nộp lại khoản thu, buộc thông báo cho chủ thể dữ liệu bị ảnh hưởng (khoản 5 Điều 53); buộc nộp lại số lợi bất hợp pháp (điểm d khoản 4 Điều 15); buộc lưu trữ dữ liệu tại Việt Nam (điểm a khoản 3 Điều 33); buộc thu hồi tên miền do thực hiện hành vi vi phạm (khoản 3 Điều 5).

\subsection{Người bán và người mua trên sàn}\label{sec:7.3}

Phạm vi đề tài tập trung vào nhà vận hành sàn (mục~\ref{sec:1.2}), nhưng Nghị định 330/2026 có điều xử phạt riêng cho người bán và người mua. Bảng dưới đây chỉ liệt kê để đối chiếu.

{\small\setstretch{1.1}
\begin{table}[H]
\centering
\caption{Một số hành vi của người bán, người mua theo Nghị định 330/2026}\label{tab:7.2}
\begin{tabularx}{\textwidth}{>{\hsize=1.4\hsize\raggedright\arraybackslash}X>{\hsize=0.8\hsize\raggedright\arraybackslash}X>{\hsize=0.8\hsize\raggedright\arraybackslash}X}
\toprule
\textbf{Hành vi} & \textbf{NĐ 330/2026} & \textbf{Cá nhân / tổ chức} \\
\midrule
Đưa lên mạng thông tin thuộc bí mật cá nhân, đời sống riêng tư (người bán) & Điểm a khoản 1 Điều 13 & 10--20 / 20--40 triệu \\
Bán, cho thuê tài khoản số có chức năng giao dịch tài chính của mình (người bán) & Khoản 1 Điều 14 & 5--10 / 10--20 triệu \\
Dùng trái phép thông tin của người khác để đăng ký tài khoản số có chức năng giao dịch tài chính & Điểm a khoản 2 Điều 14 & 10--20 / 20--40 triệu \\
Sử dụng thông tin đăng nhập, quyền truy cập tài khoản số không đúng quy định (người mua) & Điểm c khoản 3 Điều 14 & 20--30 / 40--60 triệu \\
Sử dụng tài khoản mạng xã hội của người khác không được chủ tài khoản đồng ý (người mua) & Điểm e khoản 5 Điều 14 & 50--70 / 100--140 triệu \\
\bottomrule
\end{tabularx}
\end{table}}

\subsection{Nhận xét}\label{sec:7.4}

\begin{itemize}
\item Hành vi cốt lõi là mua bán dữ liệu cá nhân không bị phạt ở các mục an ninh mạng (Mục 1 đến Mục 5) mà ở Điều 53, thuộc mục bảo vệ dữ liệu cá nhân, với mức phạt tính theo khoản thu. Đây là điều khoản duy nhất trong bảng~\ref{tab:7.1} có mức phạt tăng theo quy mô giao dịch.
\item Các nghĩa vụ mà chính sách Đăng ký, Gỡ tin, Lưu trữ ở mục~\ref{sec:5.2.2} dựa vào (xác thực tài khoản, ngăn chặn hành vi vi phạm, lưu dữ liệu trong nước) đều có điều xử phạt tương ứng ở Điều 15, 29 và 33 Nghị định 330/2026. Vì vậy, vi phạm chính sách của bộ 3P cũng là căn cứ xử lý hành chính.
\item Với vụ taphoammo.net, các hành vi đã bị khởi tố hình sự và xảy ra trước khi hai văn bản có hiệu lực. Chương này không thay đổi kết luận ở chương 3.
\end{itemize}
